\chapter*{TÓM TẮT NỘI DUNG DỰ ÁN}

Dự án trình bày việc thiết kế và phát triển hệ thống SecureChat, một giải pháp truyền thông thời gian thực (Real-time Communication) dựa trên mô hình Client-Server. Hệ thống được xây dựng bằng ngôn ngữ C++ nhằm đáp ứng hai yêu cầu cốt lõi: khả năng xử lý đồng thời (High Concurrency) và bảo mật thông tin đa lớp.

Ở phía Máy chủ (Server), kiến trúc được thiết kế theo mô hình điều khiển sự kiện bất đồng bộ (Event-Driven) thông qua cơ chế \textit{epoll} kết hợp với Thread Pool. Cấu trúc này hỗ trợ giải quyết bài toán C10K và tối ưu hóa thông lượng tại các luồng I/O mạng. Cơ sở dữ liệu SQLite được triển khai cùng cơ chế Write-Ahead Logging (WAL) nhằm duy trì tính nhất quán và tốc độ truy xuất trong môi trường đa luồng.

Về giao thức và bảo mật, hệ thống ứng dụng giao thức truyền tải nhị phân tùy chỉnh (Custom Binary Protocol) có tích hợp mã kiểm lỗi CRC32 tại Header. Cơ chế mã hóa lai (Hybrid Crypto) được áp dụng, sử dụng RSA-2048 để thiết lập khóa phiên an toàn (Handshake) và AES-256-GCM để mã hóa, đồng thời xác thực tính nguyên vẹn của toàn bộ tải trọng thông điệp (Payload).

Bên cạnh đó, dự án tích hợp các phân hệ an ninh mạng chủ động, bao gồm Hệ thống phát hiện xâm nhập (Intrusion Detector - IDS) để phân tích hành vi bất thường, Bộ giới hạn tốc độ (Rate Limiter) dựa trên thuật toán Token-Bucket nhằm phòng chống tấn công DoS/Spam, và hệ thống Nhật ký kiểm toán (Audit Logger) để lưu vết các sự kiện bảo mật. Phía Máy khách (Client) triển khai giao diện dòng lệnh (CLI), hỗ trợ đa dạng nghiệp vụ từ nhắn tin cá nhân, phòng trò chuyện nhóm, trung chuyển tệp tin ẩn danh, đến các luồng thực thi lệnh quản trị (Admin Commands) dựa trên cơ sở phân quyền (RBAC).

Kết quả nghiên cứu cung cấp một nền tảng giao tiếp hoạt động ổn định, tuân thủ các tiêu chuẩn bảo mật hiện hành và đảm bảo khả năng mở rộng kiến trúc hệ thống trong tương lai.
